\section{A conditional Frobenius boundary with quantum composition}
\label{sec:frobenius-boundary}

There is a concrete way for arithmetic Frobenius to survive at one:
retain the quantum systems carried by its moving orbits, together with
their conditional reference states. The smallest example is a qubit on
which Frobenius exchanges two distinguishable states. Its weight in the
original register vanishes, but its normalized internal physics does not.
The construction below includes independent composition and standard
subfield inclusion. It is a selected observable family; extending it to
the full arithmetic multiplication and Fourier processes remains open.
The new statements carry sketch status pending independent review.
The ``orbit checker'' cited below is
\path{theory/checks/frobenius_boundary_check.py}.

\subsection{The arithmetic comparison and its positive continuation}

For a prime $p$ and an integer $r\geq1$, let $E=\mathbb F_{p^r}$,
$H_E=\ell^2(E)$ with orthonormal basis $|x\rangle$, and
$U_E|x\rangle=|x^p\rangle$. The ordinary physical reference trace is
$\operatorname{Tr}/p^r$. We continue the prime parameter by a real variable
$t\geq1$, with $r$ fixed. This defines an algebra family; it does not posit
finite fields of noninteger cardinality.
\provenance{D1301, D1302}{the preceding arithmetic sections}
{the arithmetic Frobenius checker}{the arithmetic adjudication}

\begin{definition}[Marked Frobenius orbit algebra]
Fix $r\geq1$ and a real variable $t\geq1$, distinct from a field
characteristic. Let $\mu(n)$ be zero when a prime square divides $n$ and
otherwise $(-1)$ to the number of its prime factors. Let $\varphi(d)$ count
the integers $1\leq a\leq d$ coprime to $d$. Set
\[
 c_d(t)=\sum_{e\mid d}\mu(d/e)t^e,\qquad
 \mathcal O_r=\bigoplus_{d\mid r}M_d(\mathbb C),\qquad
 S_d|j\rangle=|j+1\bmod d\rangle,\qquad
 u_r=\bigoplus_{d\mid r}S_d.
\]
At $t=p$, choose $E=\mathbb F_{p^r}$ and an origin $b_O$ in each
Frobenius orbit $O$. Label that orbit by $\sigma_E^j(b_O)$,
$0\leq j<|O|$. The representation $R_{p,r}$ acts with the same matrix
$a_d$ on every marked orbit of length $d$. For a tower of standard
subfields in one finite field, use the same origin wherever an orbit occurs.
\end{definition}
\begin{scope}
The comparison includes the origin choices. Compatible origins for all
named embeddings or field automorphisms are not stipulated.
\end{scope}
\provenance{D1331}{definitions.md}{orbit checker B1}{not yet independently reviewed}

\begin{proposition}[Arithmetic orbit counts and matrix comparison]
\statusSketch\quad
For every prime $p$, every $r\geq1$, and each $d\mid r$, exactly $c_d(p)$
field labels have Frobenius period $d$. The marked representation is
faithful, unital and preserves adjoints, and
\[
 R_{p,r}(u_r)=U_E,\qquad
 \frac{\operatorname{Tr}(R_{p,r}(a))}{p^r}
   =\sum_{d\mid r}\frac{c_d(p)}{p^r}\operatorname{tr}_d(a_d),
 \qquad \operatorname{tr}_d=\frac{\operatorname{Tr}}d.
\]
Changing the origins gives an orbitwise unitary conjugation commuting with
$U_E$.
\end{proposition}
\begin{scope}
Marked rank-one observables can change with the origins. This representation
retains common matrices on equal-length orbits, not the full field matrix
algebra and not the commutant of Frobenius.
\end{scope}
\provenance{FRL-ORBIT}{orbit-boundary.md, Sections 1--2}
{orbit checker B1--B2}{not yet independently reviewed}

For $d\mid r$, the polynomial $X^{p^d}-X$ divides $X^{p^r}-X$ and has
derivative $-1$, so it has exactly $p^d$ roots in $E$. Decomposing these
fixed labels by their least period and applying divisor inversion gives
$c_d(p)$. Each matrix $a_d$ occurs $c_d(p)/d$ times, giving the trace
formula. This is the usual necklace polynomial multiplied by $d$; see
Trevor Hyde, \emph{Cyclotomic factors of necklace polynomials}, introduction,
p.~1. The matrix and trace continuation is the construction here.
It uses actual coordinate tests: summing $|b_O\rangle\langle b_O|$ over
length-$d$ orbit origins, then multiplying on the left and right by
$U_E^i$ and $U_E^{-j}$, produces the represented matrix unit
$|i\rangle\langle j|_d$. Frobenius and the marked origin tests therefore
generate the selected matrix algebra.

\begin{definition}[Conditional orbit boundary and reference state]
For $r\geq2$ let $e_r$ be identity on the blocks $d>1$ and zero on the
$d=1$ block, and put
\[
 \mathcal O_r^+=e_r\mathcal O_re_r,\qquad u_r^+=e_ru_re_r,
 \qquad \theta_{r,t}(a)=\sum_{d\mid r}\frac{c_d(t)}{t^r}
                     \operatorname{tr}_d(a_d),\qquad
 w_r(t)=1-t^{1-r}.
\]
The identity of $\mathcal O_r^+$ is $e_r$. For $t>1$ define
$\omega_{r,t}(a)=\theta_{r,t}(a)/w_r(t)$ on this corner. At one define
\[
 \omega_{r,1}(a)=\sum_{\substack{d\mid r\\d>1}}
               \frac{\varphi(d)}{r-1}\operatorname{tr}_d(a_d).
\]
Write $\lambda_{r,d}(t)=c_d(t)/(t^r-t)$ for $t>1$ and
$\lambda_{r,d}(1)=\varphi(d)/(r-1)$. The ordinary block-trace density is
$\sigma_{r,t}=\bigoplus_{d\mid r,d>1}\lambda_{r,d}(t)I_d/d$.
At $t=p$ the physical conditioning projection is
$P_r^{\rm mov}=\sum_{x^p\ne x}|x\rangle\langle x|$.
\end{definition}
\begin{scope}
This removes fixed basis labels. It does not remove every invariant vector:
a moving orbit still contains an invariant uniform superposition.
\end{scope}
\provenance{D1332}{definitions.md}{orbit checker B2--B3}{not yet independently reviewed}

\begin{proposition}[A faithful positive conditional endpoint]
\statusSketch\quad
For every $r\geq2$, $\theta_{r,t}$ is a faithful normalized positive
trace for $t>1$, and $\omega_{r,t}$ extends continuously to the displayed
faithful normalized positive trace on $\mathcal O_r^+$ at $t=1$.
For every integer $k$, with $\gcd(r,0)=r$,
\[
 \omega_{r,t}((u_r^+)^k)=\frac{t^{\gcd(r,k)}-t}{t^r-t}\quad(t>1),
 \qquad
 \omega_{r,1}((u_r^+)^k)=\frac{\gcd(r,k)-1}{r-1}.
\]
For any block $d>1$, the state and effect $|0\rangle\langle0|$ have
Born probability one before Frobenius and zero afterwards.
\end{proposition}
\begin{scope}
This is a positive family of selected marked orbit observables. It does
not specialize all arithmetic amplitudes. The unconditioned trace at one
has a different quotient, retaining only the scalar block.
\end{scope}
\provenance{FRL-POS}{orbit-boundary.md, Sections 2--4}
{orbit checker B2--B3}{not yet independently reviewed}

Positivity holds throughout the real interval, not just at arithmetic
points. Put $t=e^h$ with $h>0$. For $d>1$,
\[
 c_d(e^h)=\sum_{k\geq1}\frac{h^k}{k!}
 d^k\prod_{\substack{\ell\mid d\\\ell\text{ prime}}}(1-\ell^{-k}).
\]
This convergent series has positive terms. Its first coefficient is
$\varphi(d)$, so $c_d'(1)=\varphi(d)$. Since
$\sum_{d\mid r,d>1}c_d(t)=t^r-t$, division and cancellation at one give
the stated weights, which are positive and sum to one. A cyclic shift has
trace zero unless its power fixes every index, yielding the moment formula.
Thus $C([1,T],\mathcal O_r^+)$, for any finite $T>1$, is an explicit
continuous section algebra with these reference states.

For $r=2$, the entire conditional family is $M_2(\mathbb C)$ with trace
$\operatorname{tr}_2$ and Frobenius $S_2$. At $p=2$ its physical orbit is
$(\alpha,\alpha^2)$ in $\mathbb F_4$, where $\alpha^2=\alpha+1$.
The ambient probability of this qubit sector is $(t-1)/t$; the normalized
qubit remains after division by that vanishing weight. For $r=4$ the
conditional algebra is $M_2\oplus M_4$:
\[
 (\lambda_{4,2}(t),\lambda_{4,4}(t))
   =\left(\frac1{t^2+t+1},\frac{t(t+1)}{t^2+t+1}\right).
\]
Its weights change from $(1/7,6/7)$ at $t=2$ to $(1/3,2/3)$ at one.
Frobenius remains measurable on both blocks.

\subsection{Composition and retained subfield information}

\begin{definition}[Orbit processes and coherent independent composition]
Objects are finite ordered words $R=(r_1,\ldots,r_n)$, $r_j\geq2$, and
finite nonempty tagged families of such words. A word has algebra
$\mathcal O_R^+=\bigotimes_j\mathcal O_{r_j}^+$, product reference
$\omega_{R,t}=\bigotimes_j\omega_{r_j,t}$, and ambient weight
$W_R(t)=\prod_j w_{r_j}(t)$. The empty word has algebra $\mathbb C$;
tagged families have direct sum algebras. Processes are complex-linear
CPTP maps in the Schr\"odinger orientation, with ordinary sum-of-blocks
trace. Equality is equality of linear maps; composition is composition
and independent tensor pairs tags. Frobenius is conjugation by
$\bigotimes_j u_{r_j}^+$.

For $2\leq r\mid s$, let $e_{r,s}$ be identity on the blocks $d\mid r$,
$d>1$, in $\mathcal O_s^+$ and zero elsewhere. The encoding $j_{r,s}$
inserts zero blocks. The retained decoder sends $\rho$ to
\[
 \left((\rho_d)_{d\mid r,d>1},\ (1-e_{r,s})\rho(1-e_{r,s})\right)
 \quad\text{in }\mathcal O_r^+\oplus\mathcal O_s^+,
\]
with success and failure tags, respectively. Denote its reference success
probability by $h_{r,s}(t)=(t^r-t)/(t^s-t)$ for $t>1$, and by
$h_{r,s}(1)=(r-1)/(s-1)$ at one.

For $d,e\geq1$, set $g=\gcd(d,e)$ and $l=\operatorname{lcm}(d,e)$.
The reblocking is
\[
 B_{d,e}:\mathbb C^g\otimes\mathbb C^l\longrightarrow
                  \mathbb C^d\otimes\mathbb C^e,
 \qquad |a,k\rangle\longmapsto|k\bmod d,a+k\bmod e\rangle,
 \quad 0\leq a<g,\quad0\leq k<l.
\]
Its multiplicity factor is quantum. Iterated reblockings are compared
through the same ordered Cartesian basis. The word length records order
of ambient vanishing and is not a Clifford level.
\end{definition}
\begin{scope}
The CP envelope is not asserted to be the image of all arithmetic source
processes. Arithmetic comparison covers the marked observables and the
specified gates and block instruments.
\end{scope}
\provenance{D1333}{definitions.md}{orbit checker B4--B5}{not yet independently reviewed}

\begin{proposition}[Positive composition and conditional descent]
\statusSketch\quad
The word and tagged algebras above form symmetric monoidal CPTP categories
with continuous product reference states for every $t\geq1$. Continuous
CPTP matrix families evaluate functorially at one. Standard subfield
inclusions with coherent orbit origins give the stated block encodings
and retained decoders, commute with Frobenius, and compose on success
along towers. Their reference success probability is $h_{r,s}(t)$.
The reblocking is unitary and satisfies
\[
 B_{d,e}^*(S_d\otimes S_e)B_{d,e}=I_g\otimes S_l,
 \qquad M_d\otimes M_e\simeq M_g\otimes M_l.
\]
Iterated reblockings have coherent unitary connectors through the common
Cartesian basis.
\end{proposition}
\begin{scope}
Independent tensor retains coherence between simultaneous Frobenius
orbits. It is not identified with a field compositum or with a classical
sum over those orbits. Arbitrary named embeddings require their own
orbit-coordinate transport. Full decoder histories keep their tags.
The prime-field register has no moving sector and is not a conditional atom.
\end{scope}
\provenance{FRL-COMP}{orbit-composition.md, Sections 1--3}
{orbit checker B4--B5}{not yet independently reviewed}

The cycle-pair formula follows by solving two congruences for $k$, with
$a$ the residue of the difference of the two output labels modulo $g$.
Its set-level counterpart is the gcd/lcm product in the necklace algebra
(Yoshida, Section~2.2, p.~127). For two degree-two cycles, the quantum
algebra is $M_4$, and Frobenius becomes $I_2\otimes S_2$ after reblocking.
Replacing it by $M_2\oplus M_2$ dephases the two orbits
$\{00,11\}$ and $\{01,10\}$. On the state
$(|00\rangle+|01\rangle)/\sqrt2$, that dephasing reduces the return
probability from one to $1/2$.

Inclusion is equally concrete: a label belongs to $\mathbb F_{p^r}$ inside
$\mathbb F_{p^s}$ exactly when its period divides $r$. The matching
divisor blocks give the decoder. At the conditional boundary, the
degree-two inclusion into degree four succeeds on the reference state
with probability $1/3$. Two independent decoders retain all four
probabilities $1/9,2/9,2/9,4/9$.

\subsection{Higher interactions and the remaining mixed construction}

\begin{definition}[Multiplication cuts and Fourier-transported sectors]
Let $E$ be a finite field of size $Q$, and let $d\geq1$ separate controls
and a separate target carry the gate
\[
 M_E^{(d)}|x_1,\ldots,x_d,z\rangle
 =|x_1,\ldots,x_d,z+\textstyle\prod_j x_j\rangle.
\]
Put $P_{E,d}^{\rm nz}=(I-|0\rangle\langle0|)^{\otimes d}\otimes I$,
$\tau_{E,d}=\operatorname{Tr}/Q^{d+1}$, and
$v_{E,d}=P_{E,d}^{\rm nz}M_E^{(d)}P_{E,d}^{\rm nz}$.
The corner trace is
$\tau_{E,d}^{\rm nz}(a)=\tau_{E,d}(a)/\tau_{E,d}(P_{E,d}^{\rm nz})$.
On a specified word, $P_A^{\rm nz}$ tests nonzero labels precisely on
the register subset $A$.
For $E=\mathbb F_{p^r}$, let
$F_E|x\rangle=Q^{-1/2}\sum_y\exp(-2\pi i\operatorname{Tr}_{E/\mathbb F_p}(xy)/p)|y\rangle$
and put $P_r^{\rm dual}=F_EP_r^{\rm mov}F_E^*$.
\end{definition}
\begin{scope}
Nonzero controls, moving Frobenius labels and their Fourier-transported
cuts are distinct tests.
\end{scope}
\provenance{D1306, D1308, D1334}{definitions.md}{orbit checker B6}
{arithmetic definitions admitted; new cuts not yet independently reviewed}

\begin{proposition}[Active weights and interaction order]
\statusSketch\quad
For every finite field $E$, every $d\geq1$, and the traces just specified,
write $w=((Q-1)/Q)^d$. Then
\[
 \tau_{E,d}(P_{E,d}^{\rm nz})=w,\quad
 \tau_{E,d}(M_E^{(d)})=1-w,\quad
 \tau_{E,d}((M_E^{(d)}-I)^*(M_E^{(d)}-I))=2w.
\]
The corner gate is unitary, with trace zero and conditional squared
difference from its corner identity equal to two. Frobenius and field
inclusions preserve these nonzero-control cuts. On one word,
$P_A^{\rm nz}P_B^{\rm nz}=P_{A\cup B}^{\rm nz}$, with reference weight
$((Q-1)/Q)^{|A\cup B|}$. Independently conditioned moving-orbit words obey
\[
 W_R(t)=\left(\prod_{j=1}^n(r_j-1)\right)(t-1)^n
         +O((t-1)^{n+1}).
\]
\end{proposition}
\begin{scope}
These are trace and counting identities, not operator-norm convergence.
The multiplication family's order $d$ matches its interaction arity and
exact Clifford level $d+1$; no general equivalence is asserted.
\end{scope}
\provenance{FRL-ACTIVE}{orbit-composition.md, Sections 4--5}
{orbit checker B6}{not yet independently reviewed}

There are $(Q-1)^d$ active control tuples. Each gives a nonzero translation
of the target, with no fixed labels; all other tuples give the identity.
This proves the traces by counting permutation fixed points. Overlapping
controls are counted once, which explains why their vanishing orders need
not add. Independent moving factors do multiply, and must retain that
order information even after their individual normalizations.

The extension to mixed arithmetic processes needs more structure. In three
$\mathbb F_4$ registers, multiplication sends
$(\alpha,\alpha,\alpha^2)$ to $(\alpha,\alpha,0)$, leaving the tensor
of moving-label sectors. Also, for the normalized trace-fibre inclusion
$V|a\rangle=|\ker\operatorname{Tr}_{E/K}|^{-1/2}
\sum_{\operatorname{Tr}_{E/K}(x)=a}|x\rangle$, the map
$\mathbb F_2\subset\mathbb F_4$ has
$V|0\rangle=(|0\rangle+|1\rangle)/\sqrt2$; its nonzero cut does not
intertwine with $V$. Finally the Fourier matrix has constant squared
entry modulus $1/Q$, giving
\[
 \frac{\operatorname{Tr}(P_r^{\rm mov}P_r^{\rm dual})}{p^r}
 =(1-p^{1-r})^2.
\]
This is an overlap of possibly noncommuting projections, not an intersection
dimension. It supplies a definite second-order mixed datum. The next
construction must retain the two charts, their measured transitions and
the arithmetic multiplication and transfer relations in one positive
section algebra. The present Frobenius orbit family supplies a concrete
positive starting point for that task.
\provenance{D1304, FRL-ACTIVE}{orbit-composition.md, Section 5}
{orbit checker B6}{supporting scope calculations; no additional promotion}
